\section{Budget système : mémoire, latence, énergie, calcul}
\label{sec:system}

Un gain de quantification ne se juge pas sur les poids seuls. C'est ici que se tranche l'argument
tripartite rappelé en introduction --- moins de mémoire, moins de temps, moins d'énergie : cette section
rassemble les quatre budgets d'une même tête EWC déployée, tous mesurés sur la même carte, et prend chaque
promesse séparément. Les trois instrumentations correspondantes (composantes de la RAM, segmentation du
noyau au compteur de cycles, estimateurs d'énergie) sont décrites en \S\ref{sec:method}.

\subsection{La RAM des poids n'est pas la RAM du système}

L'INT8 divise bien par~4 la RAM \emph{des poids} (\num{2688} à \num{672}~octets sur Monitoring,
\num{2816} à \num{704} sur Pronostia). Ce décompte porte sur les matrices de poids seules, celles que la
quantification touche ; les biais restent flottants et exacts (\S\ref{sec:method}), si bien que le nombre
de paramètres de la tête, \num{722} et \num{754}, est légèrement supérieur. Mais la RAM réellement consommée par le système est
$\texttt{.data} + \texttt{.bss} + \text{pic de pile}$, et cette somme ne bouge pratiquement pas : elle passe
de \num{105300} à \num{105324}~octets sur Monitoring, soit un rapport INT8/FP32 de \num{1.0002}.
Le détail par composante est en Annexe~\ref{app:budget} (Tableau~\ref{tab:annex-ram}). Les deux
configurations occupent \num{40.17}\,\% et \num{40.18}\,\% du budget de \SI{256}{\kilo\byte}
(\num{42.03}\,\% et \num{42.04}\,\% sur Pronostia).

L'explication est simple et vaut d'être dite clairement : à cette échelle de modèle, les poids ne sont pas
le poste dominant --- les tampons de flux, la pile et le reste du firmware le sont. Annoncer un « $\div 4$
de RAM » sans préciser \emph{de quelle} RAM on parle serait donc trompeur. Le $\div 4$ est réel et utile,
mais il porte sur un poste minoritaire de ce système-ci ; il deviendrait décisif sur un modèle où les poids
dominent.

\begin{figure}[t]
\centering
\includegraphics[width=0.85\linewidth]{fig8_ram_total.png}
\caption{RAM totale \emph{mesurée sur carte} : \texttt{.data} $+$ \texttt{.bss} $+$ pic de pile, en
regard du budget SRAM. Le total INT8 égale le total FP32, alors même que les poids sont divisés par~4.}
\label{fig:ram}
\end{figure}

\subsection{Où passe la latence INT8}

L'instrumentation cycle-à-cycle du kernel INT8 chiffre le paradoxe annoncé. Sur Monitoring, le MAC entier
coûte \num{6841}~cycles, la requantification \num{3760} et la déquantification \num{200} (respectivement
\num{6729}, \num{3794.5} et \num{200} sur Pronostia). Le total atteint \SI{74}{\micro\second} contre
\SI{48}{\micro\second} en FP32 sur Monitoring et \SI{50}{\micro\second} sur Pronostia.

Le poste décisif n'est donc pas la multiplication : c'est la \emph{requantification}, un tiers du budget, et
elle n'existe tout simplement pas dans le chemin FP32. Le MAC entier, lui, n'apporte aucun gain face au FPU
matériel du Cortex-M4. Le surcoût INT8 n'est pas une anomalie de mise en \oe uvre mais la conséquence
arithmétique d'un format entier exécuté sur une cible qui calcule nativement en flottant.

\subsection{Énergie : trois estimateurs, jamais fusionnés}

L'énergie a été mesurée à la sonde X-NUCLEO-LPM01A par trois méthodes indépendantes, que nous reportons
\emph{côte à côte} sans jamais les moyenner --- leur désaccord est lui-même un résultat. La régression du
courant en fonction de la cadence donne \SI{67.65}{\micro\joule} par inférence en FP32 et
\SI{67.47}{\micro\joule} en INT8 ; le delta contre un repos \texttt{WFI} donne \SI{146.77}{\micro\joule} et
\SI{145.03}{\micro\joule} ; le lot d'inférences par trame, qui isole le calcul du coût de la trame UART,
donne \SI{5.08}{\micro\joule}. Le Tableau~\ref{tab:annex-energy} les range côte à côte.
L'ordre de grandeur commun aux deux premières et l'écart d'un facteur~10 avec
la troisième délimitent exactement ce que chacune mesure : les deux premières intègrent le coût de
communication, la dernière non.

Sur la question qui nous occupe, les trois convergent : \textbf{l'INT8 ne réduit pas la consommation}. L'écart
mesuré est de \SI{0.18}{\micro\joule} pour une incertitude combinée de \SI{2.34}{\micro\joule}, soit très en deçà du
seuil de signification. Le paradoxe de la latence se double donc d'un paradoxe d'énergie, et la
justification de l'INT8 se réduit à la mémoire.

Le balayage de fréquence ouvre en revanche un levier réel. L'énergie par inférence croît
monotonement avec l'horloge --- \SI{170.05}{\micro\joule} à \SI{45}{\mega\hertz},
\SI{192.64}{\micro\joule} à \SI{90}{\mega\hertz}, \SI{214.61}{\micro\joule} à
\SI{180}{\mega\hertz} --- pendant que la latence reste très en deçà du budget Gap~2, jusqu'à
\SI{198}{\micro\second} au point le plus lent. Autrement dit, \emph{la marge de latence est convertible
en autonomie} : au régime le plus sobre mesuré, une pile de \SI{2000}{\milli\ampere\hour} tient
\SI{72.95}{\hour} à une inférence par seconde. Deux cellules restent « à mesurer » : l'énergie par mise à
jour continue, dont les régressions ne sont pas publiables en l'état, et le profil par phase.

\begin{figure}[t]
\centering
\includegraphics[width=0.9\linewidth]{fig9_energy_estimators.png}
\caption{Énergie \emph{mesurée sur carte}. À gauche, les trois estimateurs côte à côte ; à droite,
l'arbitrage fréquence, où la marge de latence s'échange contre de l'énergie.}
\label{fig:energy}
\end{figure}

\subsection{Coût de calcul}

Analytiquement, la tête compte \num{722} paramètres sur Monitoring (\num{754} sur Pronostia), pour
\num{672} MAC et \num{1344} \gls{FLOP} par inférence (\num{704} et \num{1408}). Le compte en \emph{opérations
binaires} donne le ratio théorique attendu : \num{1376256} BOPs en FP32 contre \num{86016} en INT8, soit un
facteur~16 --- c'est-à-dire $(32/8)^2$.

Ce facteur~16 est précisément ce que la carte \emph{ne} restitue pas. Confronter les deux colonnes de la
Figure~\ref{fig:cost} est la manière la plus directe de voir pourquoi une métrique de coût analytique ne
peut pas remplacer une mesure : le \gls{BOP} compte des bits d'opérande, pas des cycles d'une unité flottante
existante.

\begin{figure}[t]
\centering
\includegraphics[width=0.9\linewidth]{fig10_compute_cost.png}
\caption{Coût de calcul \emph{théorique} (gauche, hachuré) contre latence \emph{mesurée sur carte}
(droite). Le gain de \num{16} en BOPs ne se traduit par aucun gain de temps.}
\label{fig:cost}
\end{figure}
